% year: תשפ"ו
% semester: א
% moed: ב
% number: 4ב
% points: 10
% categories: ode
% official_solution: yes
% source: exams/מבחנים/תשפו/sol B.pdf

\begin{question}
מצאו את הפתרון הכללי של המשוואה הדיפרנציאלית
\[ \sqrt{1 - x^2} \cdot y' = y^2 \arcsin x \]
\end{question}

\begin{hint}
המשוואה פרידה. אחרי ההפרדה, באגף של $x$ הציבו $t = \arcsin x$.
\end{hint}

\begin{solution}
המשוואה מוגדרת עבור $-1 < x < 1$, ושם $\sqrt{1 - x^2} > 0$. נכתוב אותה בצורה
\[ \frac{dy}{dx} = y^2 \cdot \frac{\arcsin x}{\sqrt{1 - x^2}}, \]
וזו משוואה פרידה.

\textbf{הפתרון הקבוע:} $y \equiv 0$ מקיים את המשוואה ($0 = 0$), ולכן הוא פתרון.

\textbf{עבור $y \ne 0$:} נפריד משתנים:
\[ \int \frac{dy}{y^2} = \int \frac{\arcsin x}{\sqrt{1 - x^2}}\,dx. \]
באגף שמאל מתקבל $-\frac{1}{y}$. באגף ימין נציב $t = \arcsin x$, $dt = \frac{dx}{\sqrt{1 - x^2}}$:
\[ \int \frac{\arcsin x}{\sqrt{1 - x^2}}\,dx = \int t\,dt = \frac{t^2}{2} + C = \frac{\arcsin^2 x}{2} + C. \]
לכן
\[ -\frac{1}{y} = \frac{\arcsin^2 x}{2} + C \quad\Longrightarrow\quad y = \frac{-1}{\frac{\arcsin^2 x}{2} + C} = \frac{-2}{\arcsin^2 x + 2C}. \]

\textbf{תשובה:} הפתרון הכללי הוא
\[ \boxed{y = \frac{-2}{\arcsin^2 x + C}}, \quad C \in \R \]
(בכל קטע שבו המכנה אינו מתאפס), ובנוסף הפתרון $y \equiv 0$.
\end{solution}
